% ============================================================
%   Chapter 2  --  State of the Art and Theoretical Framework
% ============================================================

\chapter{State of the Art and Theoretical Framework}\label{ch:sota}

This chapter reviews the concepts that underpin the rest of the document. It first formalises the
Reinforcement Learning problem (Section~\ref{sec:sota-rl}), then introduces Proximal Policy
Optimization and Generalised Advantage Estimation
(Sections~\ref{sec:sota-ppo}--\ref{sec:sota-gae}), surveys the Safe RL literature with a focus on
shielding (Sections~\ref{sec:sota-saferl}--\ref{sec:sota-shielding}), introduces the kinematic
bicycle model used by the adaptive shield (Section~\ref{sec:sota-bicycle}), and finishes with an
overview of the CARLA simulator and related work
(Sections~\ref{sec:sota-carla}--\ref{sec:sota-related}).

\section{Reinforcement Learning fundamentals}\label{sec:sota-rl}

The standard formulation of Reinforcement Learning relies on Markov Decision Processes
(MDPs)~\cite{sutton2018rl,puterman1994mdp}. A (discrete-time, discounted) MDP is a tuple
\(\mathcal{M} = (\mathcal{S}, \mathcal{A}, P, r, \gamma, \rho_0)\) where \(\mathcal{S}\) is the
state space, \(\mathcal{A}\) is the action space, \(P(s'\mid s,a)\) is the transition kernel,
\(r(s,a)\) is the reward function, \(\gamma \in [0,1)\) is the discount factor and \(\rho_0\)
is the initial-state distribution.

A stochastic policy \(\pi_\theta: \mathcal{S}\to\Delta(\mathcal{A})\), parameterised by
\(\theta\), induces a distribution over trajectories \(\tau = (s_0, a_0, s_1, a_1, \ldots)\).
The \emph{return} is the discounted cumulative reward:
\begin{equation}
  G(\tau) \;=\; \sum_{t=0}^{\infty} \gamma^t\, r(s_t, a_t).
  \label{eq:return}
\end{equation}
The training objective is to maximise the expected return:
\begin{equation}
  J(\theta) \;=\; \mathbb{E}_{\tau \sim \pi_\theta}\!\left[ G(\tau) \right].
  \label{eq:objective}
\end{equation}

Two central quantities guide optimisation. The \emph{state-value function} \(V^{\pi_\theta}(s)\)
measures the expected discounted return starting from state \(s\) under policy \(\pi_\theta\).
The \emph{action-value function} \(Q^{\pi_\theta}(s,a)\) additionally conditions on taking
action \(a\). The \emph{advantage} \(A^{\pi_\theta}(s,a) = Q^{\pi_\theta}(s,a) - V^{\pi_\theta}(s)\)
measures how much better than average a given action is at state \(s\).

\subsection{Policy gradient theorem}\label{sec:sota-pg}

Policy gradient methods directly optimise Equation~\eqref{eq:objective} by gradient ascent on
\(\theta\). The policy gradient theorem~\cite{sutton1999policygradient} provides a tractable
expression for the gradient:
\begin{equation}
  \nabla_\theta J(\theta) \;=\; \mathbb{E}_{\tau \sim \pi_\theta}\!\left[
     \sum_{t=0}^{\infty} \nabla_\theta \log \pi_\theta(a_t \mid s_t)\; A^{\pi_\theta}(s_t, a_t)
  \right].
  \label{eq:pg}
\end{equation}
In practice the advantage is replaced by a sample-based estimator \(\hat{A}_t\). Modern on-policy
methods such as A2C~\cite{mnih2016asynchronous}, TRPO~\cite{schulman2015trpo} and PPO~\cite{schulman2017ppo} all build on
Equation~\eqref{eq:pg}, differing mainly in how they constrain the update step size.

\section{Proximal Policy Optimization}\label{sec:sota-ppo}

Vanilla policy gradient is sensitive to the step size: a single large update can move the policy
so far from the data distribution that subsequent rollouts are unreliable. TRPO~\cite{schulman2015trpo}
resolves this by constraining the Kullback-Leibler divergence between old and new policy at each
update, but doing so requires second-order information. PPO~\cite{schulman2017ppo} achieves a
similar trust-region effect with a simple first-order surrogate.

Let \(r_t(\theta) = \pi_\theta(a_t \mid s_t)\,/\,\pi_{\theta_{\text{old}}}(a_t \mid s_t)\)
be the importance-sampling ratio. The \emph{clipped surrogate objective} is
\begin{equation}
  \mathcal{L}^{\mathrm{CLIP}}(\theta)
   \;=\; \mathbb{E}_t\!\left[
     \min\!\Big(
        r_t(\theta)\,\hat{A}_t,\;
        \mathrm{clip}\!\big(r_t(\theta),\, 1-\epsilon,\, 1+\epsilon\big)\,\hat{A}_t
     \Big)
   \right],
  \label{eq:ppo-clip}
\end{equation}
where \(\epsilon = 0.2\). As Figure~\ref{fig:ppo-clip} illustrates, the clipping prevents the
ratio \(r_t(\theta)\) from drifting far from one: for a positive advantage the policy cannot
over-exploit a good action beyond \(1+\epsilon\), and for a negative advantage it cannot
over-penalise beyond \(1-\epsilon\). This keeps updates within an approximate trust region
without ever computing a Hessian.

The full PPO loss also includes a squared-error value-function regression term weighted by
\(c_V = 0.25\) and an entropy bonus \(c_H \cdot \mathcal{H}[\pi_\theta]\) to encourage
exploration.

\begin{figure}[H]
  \centering
  \includegraphics[width=0.85\textwidth]{imgs/ppo_surrogate}
  \caption{PPO clipped surrogate objective. For a positive advantage (solid) the ratio is
           capped at \(1+\epsilon=1.2\), preventing over-exploitation of a good action. For a
           negative advantage (dashed) it is capped at \(1-\epsilon=0.8\). Both caps limit how
           far the new policy can deviate from the data-collecting policy.}
  \label{fig:ppo-clip}
\end{figure}

\subsection{Stabilisers and heuristics}\label{sec:sota-ppo-stabilisers}

Empirical studies~\cite{andrychowicz2020matters,engstrom2020implementation} have shown that
several implementation details have a first-order effect on PPO's sample efficiency and
stability. The ones relevant to this project are:

\begin{description}
  \item[\textbf{Value clipping.}] The value-function update is clipped so the critic cannot make
        overly aggressive updates within a single optimisation batch. The clipping magnitude
        \(\xi\) mirrors the PPO ratio clipping.
  \item[\textbf{KL early stopping.}] An unbiased estimator~\cite{schulman2020approxkl} of the
        KL divergence between old and new policy is computed after every minibatch. If the
        rolling average exceeds a target \(D^\star_{\mathrm{KL}} = 0.08\), the remaining epochs
        of the current update are skipped to prevent destructive updates.
  \item[\textbf{Observation normalisation.}] The 19-dimensional lane/marking/vehicle/route
        block is rescaled online to zero mean and unit variance using a running estimator. The
        rationale for normalising only the vector block (not the 720 LIDAR dimensions) is
        discussed in Section~\ref{sec:design-obsnorm}.
  \item[\textbf{Cosine learning-rate annealing.}] The learning rate decays from its initial
        value \(\eta_{\max}\) to a minimum \(\eta_{\min}\) following a cosine curve over the
        expected total number of PPO updates.
\end{description}

\section{Generalised Advantage Estimation}\label{sec:sota-gae}

The advantage \(A^{\pi_\theta}(s_t, a_t)\) in Equation~\eqref{eq:pg} must be estimated from
sampled rollouts. A plain Monte Carlo estimate has zero bias but high variance, whereas a one-step
TD residual \(\delta_t^V = r_t + \gamma V_\theta(s_{t+1}) - V_\theta(s_t)\) has low variance
but high bias. Generalised Advantage Estimation (GAE)~\cite{schulman2016gae} interpolates between
the two via an exponentially-weighted sum of TD residuals:
\begin{equation}
  \hat{A}_t^{\mathrm{GAE}(\gamma,\lambda)}
    \;=\; \sum_{\ell=0}^{\infty} (\gamma\lambda)^\ell\, \delta_{t+\ell}^V,
  \label{eq:gae}
\end{equation}
where \(\lambda \in [0,1]\) is the GAE parameter and \(\delta_{t+\ell}^V\) is the one-step TD
residual evaluated \(\ell\) steps ahead. Figure~\ref{fig:gae-lambda} shows how \(\lambda\) controls
the weight \((\gamma\lambda)^\ell\) placed on each residual, and therefore the bias-variance
trade-off: \(\lambda = 0\) keeps only the first residual (the one-step TD advantage: low variance,
biased), \(\lambda = 1\) weights all residuals by \(\gamma^\ell\) alone (the Monte Carlo advantage:
unbiased, high variance), and the value \(\lambda = 0.95\) used in this work sits between the two.
When an episode terminates due to a collision or going off-road, the critic bootstrap value is set to zero. When an episode is truncated by a step-count timeout, the critic value at the final observed state is
used to approximate the missing return tail.

\begin{figure}[H]
  \centering
  \includegraphics[width=0.85\textwidth]{imgs/gae_lambda}
  \caption{Generalised Advantage Estimation weights the one-step TD residual
           \(\delta_{t+\ell}^V\) by \((\gamma\lambda)^\ell\) (here \(\gamma = 0.99\)). The
           parameter \(\lambda\) interpolates between the low-variance, biased one-step TD estimate
           (\(\lambda = 0\), only \(\ell = 0\) contributes) and the unbiased, high-variance Monte
           Carlo estimate (\(\lambda = 1\)). The choice \(\lambda = 0.95\) gives a practical
           balance.}
  \label{fig:gae-lambda}
\end{figure}

\section{Safe Reinforcement Learning}\label{sec:sota-saferl}

The comprehensive survey by Garc\'ia and Fern\'andez~\cite{garcia2015comprehensive} organises
Safe RL into two broad categories: \emph{risk-sensitive} approaches that modify the optimality
criterion, and \emph{exploration-modifying} approaches that constrain the set of actions available
during learning. Four families are particularly relevant to autonomous driving:

\begin{itemize}
  \item \textbf{Constrained MDPs and CPO.} Safety is expressed as one or more auxiliary cost
        functions, and the policy is trained to maximise reward subject to cost-rate constraints.
        Achiam et al.~\cite{achiam2017cpo} introduced Constrained Policy Optimization (CPO), a
        TRPO-style algorithm that provides approximately monotonic constraint satisfaction.
  \item \textbf{Lagrangian and primal-dual methods.} A more scalable approach forms the dual
        problem with an adaptive multiplier for each constraint and solves it via alternating
        updates on the policy and the multiplier~\cite{stooke2020responsive}. Variants underpin
        much of modern Safe RL.
  \item \textbf{Safety-aware reward shaping.} Penalty terms in the reward function discourage
        unsafe behaviour. This is lightweight but provides no hard guarantees: the agent can
        still attempt unsafe actions during exploration. The reward shaper of
        Section~\ref{sec:design-reward} belongs to this family.
  \item \textbf{Runtime shielding.} A verifier inspects each proposed action and substitutes a
        safe fall-back if a safety specification would be violated. Shielding is orthogonal to
        the learning algorithm and can be composed with any RL backbone. Alshiekh et
        al.~\cite{alshiekh2018shielding} formalise the approach for finite systems with
        temporal-logic specifications, and Jansen et al.~\cite{jansen2020shielded} extend it to
        probabilistic settings. The shields in this project reason about predicted trajectories
        rather than formal logic.
\end{itemize}

\section{Safety shields}\label{sec:sota-shielding}

A runtime safety shield can be formalised as a function
\begin{equation}
  \Sigma: \mathcal{S}\times\mathcal{A} \to \mathcal{A},
  \qquad
  a^{\star}_t \;=\; \Sigma(s_t,\, a_t),
  \label{eq:shield-fn}
\end{equation}
that maps the current state and the agent's proposed action to a safe executed action. Two design
principles recur throughout the literature:

\begin{description}
  \item[\textbf{Minimal interference.}] The shield should override the agent's action only when
        strictly necessary, and the override should be as close as possible to the proposed
        action. In practice this is implemented by projecting onto a set of certified-safe
        actions, or by searching through a priority-ordered list of candidates.
  \item[\textbf{Credit assignment through the executed action.}] When the shield intervenes,
        the environment sees \(a^\star_t\), not the policy's \(a_t\). Storing the proposed
        action in the replay buffer would push the policy towards an action that was never
        executed. The framework addresses this by always storing \(a^\star_t\) and
        re-evaluating its log probability under the current policy
        (Section~\ref{sec:impl-shield-credit}).
\end{description}

Both shields in this project are \emph{post-posed}: they are inserted between the policy and the
simulator as a \texttt{gymnasium.Wrapper} and inspect the output of the policy after it has been
sampled.

\section{Kinematic bicycle model}\label{sec:sota-bicycle}

The adaptive shield must predict the vehicle's future position to verify whether a candidate
action will keep it within the lane. It uses a \emph{kinematic bicycle model}~\cite{rajamani2012vehicle,kong2015kinematic}, a well-established approximation that collapses each axle
into a single virtual wheel and tracks the rear axle position \((x,y)\), heading \(\theta\) and
longitudinal speed \(v\) (Figure~\ref{fig:bicycle-model}). The front-wheel steering angle
\(\delta\) and a longitudinal acceleration command drive the state forward in time steps of
\(\Delta t = 0.05\,\mathrm{s}\).

\begin{figure}[H]
  \centering
  \includegraphics[width=0.65\textwidth]{imgs/bicycle_model}
  \caption{Kinematic bicycle model geometry used by the adaptive shield. The rear axle is at
           \((x,y)\) with heading \(\theta\). Control inputs are the front-wheel steering
           \(\delta\) and a longitudinal acceleration. The model advances one step
           \(\Delta t = 0.05\,\mathrm{s}\) using a closed-form arc integration to avoid
           Euler-step drift at large steering angles.}
  \label{fig:bicycle-model}
\end{figure}

For small \(|\delta|\) the state advances with a simple forward-Euler step. For non-negligible
\(|\delta|\), the exact arc geometry is used to avoid integration drift. The discrete equations
are implemented in \texttt{src/Adaptative\_Shield/BicycleModel.py}. Kong et
al.~\cite{kong2015kinematic} show that, at the speeds and curvatures typical of urban driving,
the kinematic model closely approximates the full dynamic bicycle model for planning horizons
below one second, which covers the adaptive shield's longest horizon of ten steps
(\(0.5\,\mathrm{s}\)). The vehicle parameters for CARLA's Tesla Model~3 are in
Table~\ref{tab:bicycle-params} (Chapter~\ref{ch:design}).

\section{CARLA simulator}\label{sec:sota-carla}

CARLA~\cite{dosovitskiy2017carla} is an open-source simulator for autonomous driving research
built on Unreal Engine 4. Its features relevant to this thesis are summarised below.

\begin{itemize}
  \item \textbf{Client\textendash server architecture.} The simulator runs as a server process
        (\texttt{CarlaUE4.exe}) that holds the world state and renders the scene, while the
        Python client connects to it over a TCP port. This decoupling makes it straightforward
        to run training and the simulator on different machines.
  \item \textbf{Synchronous mode.} By default CARLA runs asynchronously. When
        \texttt{synchronous=True} the server waits for a \texttt{world.tick()} call from the
        client before advancing one fixed time step (\(\Delta t = 0.05\,\mathrm{s}\), i.e.\
        \(20\,\mathrm{Hz}\)). This is a prerequisite for reproducible experiments.
  \item \textbf{Waypoint API.} The road network is represented as an OpenDRIVE map from which
        CARLA exposes a lattice of \emph{waypoints} along lane centres. Each waypoint carries
        lane width, lane type, road curvature, lane-marking type and the allowed lane-change
        directions. The Waypoint API is the primary source of geometric information for both the
        reward shaper and the safety shields in this project.
  \item \textbf{Sensors.} CARLA provides a library of virtual sensors. The framework uses three:
        a semantic LIDAR, a collision detector and a lane-invasion detector. The semantic LIDAR
        tags every ray hit with the \texttt{carla.CityObjectLabel} of the struck primitive,
        enabling the three-channel decomposition of Section~\ref{sec:design-obs-lidar}.
  \item \textbf{Traffic Manager.} Non-playable vehicles (NPCs) are governed by CARLA's internal
        Traffic Manager, which follows the simulator's own traffic rules. The framework uses it
        to instantiate the NPC density required by each curriculum stage.
  \item \textbf{Predefined towns.} Town04 is a highway loop with multi-lane segments and
        entrance ramps, and it is the main benchmarking map used in
        Chapter~\ref{ch:experimentation}. Towns 01--03 cover urban scenarios, and Town05 offers
        maximum difficulty.
\end{itemize}

\section{Related work}\label{sec:sota-related}

\paragraph{Deep RL for autonomous driving.}
The surveys by Kiran et al.~\cite{kiran2022drlsurvey} and Aradi~\cite{aradi2022survey} offer
broad panoramas of the field. PPO is the most common backbone for continuous-control experiments
on CARLA, with notable examples including Chen et al.~\cite{chen2019carlarl} on urban driving and
CARLA Leaderboard submissions combining PPO with learned perception~\cite{chitta2023transfuser}.

\paragraph{Safety shields in simulated driving.}
Shielding has been applied to benchmarks such as HighwayEnv~\cite{leurent2018highway} and CommonRoad~\cite{althoff2017commonroad}. Krasowski et
al.~\cite{krasowski2020safe} use reachability analysis to certify lane-change manoeuvres,
Wabersich and Zeilinger~\cite{wabersich2021filter} combine learning-based control with a
model-predictive safety filter, and MetaDrive~\cite{li2023metadrive} provides a lighter-weight
shielding interface that inspired the early prototypes of this project. The present work
contributes a \emph{pair} of interchangeable shields in the CARLA ecosystem.

\paragraph{Curriculum learning.}
The general idea of presenting tasks in order of increasing difficulty dates back to Bengio et
al.~\cite{bengio2009curriculum}. Applications to RL range from procedural domain
randomisation~\cite{akkaya2019solving} to self-play curricula. The curriculum manager of
Section~\ref{sec:design-curriculum} includes a \emph{rollback} mechanism driven by empirical
instability on the 0--20~NPC progression, which is uncommon in the curriculum-learning
literature.

\paragraph{Bicycle model for trajectory prediction.}
Kong et al.~\cite{kong2015kinematic} show that, at speeds and curvatures typical of urban
driving, the kinematic model closely approximates the dynamic bicycle model for sub-second
horizons. This justifies its use inside the adaptive shield
(Section~\ref{sec:sota-bicycle}).

\clearpage
